\scanpage{24}
% lines 1–4
\uline{deux} mbs de la $1^{\mathrm{ère}}$ sont multipliés alors par $(\mu\lambda)^2=\mu^2\lambda^2$, et les \uline{deux} mbs de la seconde par $(\mu^2\lambda)^2=\mu^4\lambda^2$. On vérifie ensuite les formules dans le cas particulier (21).

% lines 4–11
On peut se simplifier le travail en notant que pour l'une et l'autre formule, les deux mbs sont des formes invariantes sous l'action du groupe \inserted{spécial} orthogonal $\underline{\mathrm O}(q,\omega)$ (sur $\operatorname{Spec}\unsure{Z|\mathbb Z|\mathbf Z}$ si on veut) or il n'y a (en car zéro, si on y tient -- celà suffit -- qu' une seule forme invariante sym, à scalaire près, ce qui montre la proportionnalité.

% lines 11–16
Pour avoir le coefficient correctif, il suffit de calculer la valeur des 2 mbs pour un couple d'arguments pour lesquels ils soient non nuls, p.ex. $x'=y'=H$, $x=y=\widetilde H$. On trouve

% lines 17–22
\[
\begin{aligned}
\operatorname{Cas}_{E'}(H,H)&=1+1+0=2\\
\varphi'(H,H)&=1\\
\operatorname{Cas}_E(\widetilde H,\widetilde H)
&=\operatorname{Cas}_{E'}(2H,2H)=4\operatorname{Cas}_{E'}(H,H)=8\\
&\quad(\text{ou encore }=4+4+0)\\
\varphi(\widetilde H,\widetilde H)&=2\\
\Delta&=-2
\end{aligned}
\]

% lines 23–26
On note qu'en vertu du \deleted{\unsure{bon|beau|bref}} principe de conservation des identités polynômiales, il suffit de prouver (74) sur un corps de car. 0 alg. clos, \deleted{mais \unsure{alors} on \unsure{peut}} \inserted{avec $f$ un iso i.e.} $\varphi$ associée

% lines foot
% archive sheet number: 23
